\scanpage{85}
tel que la décomposition de
$P_n(X_1\times\cdots\times X_r;Y)$ en des $P_{\nu_1\ldots\nu_r}(X_1,\ldots,X_r;Y)$
respecte les op. de $k$, i.e. $f\in P_{\nu_1\ldots\nu_r}$
implique $f\circ(\mathrm{id}_{\lambda_1}\times\cdots\times\mathrm{id}_{\lambda_r})=\lambda_1^{\nu_1}\cdots\lambda_r^{\nu_r}f$
pour tt $(\lambda_1,\ldots,\lambda_r)\in k^r$.

\medskip\noindent\hfil\shortstack[r]{\rule{3cm}{0.4pt}\\[-0.9ex]\rule{0.9cm}{0.4pt}}\hfil\par

\uline{Ex} La catégorie des $k$-modules
est une $\otimes$-pd-catégorie
De $=$ si $\unsure{\mathscr X|\mathscr E}$ un topos avec une
$k$-\deleted{\unsure{Alg.}}\unsure{Algèbre} $K$, la cat. des $K$-Modules
\deleted{et on peut} [ De $=$ on peut prendre des
$K$-Modules gradués, \deleted{\unsure{ent|par}} \inserted{par un groupe}
quelconque, \deleted{\unsure{laquelle} \unsure{soit}} on considère des
\margin{?}
produits $\Gamma^n(M^i)$ pour les degrés $ni$. ]

\medskip\noindent\hfil\shortstack[r]{\rule{3cm}{0.4pt}\\[-0.9ex]\rule{0.9cm}{0.4pt}}\hfil\par

Il y a diverses questions dans le cas
``abstrait''. Par exemple : quand
et comment \unsure{peuton|peut-on} \inserted{($\forall M\in\operatorname{Ob}P_1$)} sur
$\Gamma^*(M)$ une structure ``à p.d.'',
i.e. des hom de degré $1$
\[
\Gamma^p(\Gamma^q(M))\unsure{-|\longrightarrow}\Gamma^{pq}(M)
\]
satisfaisant aux axiomes des p.d. ?
% archive sheet number: 84
